\section{The anomalous \texorpdfstring{$\Z_2$}{Z2} symmetry of the Levin--Gu edge}%
\label{sec:asymex_czy}

The edge Hamiltonian of the Levin--Gu model, unitarily equivalent to the
XX chain, commutes with the $\Z_2$ symmetry
$U_N=\prod_nCZ_{n,n+1}\prod_nZ_nX_n$ on a ring of $N$
qubits~\cite[Section~II]{Liu2025Bialgebraic}. The source writes it as a
bond-two matrix product operator, reduces the product of two copies to a
bond-three tensor $A_0$ padded by a zero block, and states that $A_0$ is
not injective, cannot be further decomposed into a direct sum of
injective matrix product operators, and generates the identity on
periodic chains. This section reads these statements through the
multi-block asymmetric compression theorem
(Theorem~\ref{thm:asym_multiblock}). The associator of the fusion
tensors, $\omega(1,1,1)=-1$ in~\cite[Section~II]{Liu2025Bialgebraic},
is proved below. The unrestricted scalar relation on the full virtual
space is corrected to a relation on the physical images, as explained
in~\cite{gap:lmsvkl25_czy_associator}.

The first physical index is the output and the second the input, as in
the source. For a word $w$ in the index pairs, $B^w$ denotes the product
of the matrices of $B$ along $w$, as in Section~\ref{sec:asymex_czx}.

\begin{definition}[The CZY tensor and its stacked square]
    \label{def:asymex_czy_tensors}
    \lean{CZYCompression.czyTensor, CZYCompression.czySquare,
        CZYCompression.czySquareTarget}
    \leanok
    \uses{def:mpo_tensor, def:mpo_to_mps, def:mpdo_operator_product}
    The CZY tensor $A_1$ has bond dimension two and the only nonzero
    matrices
    \begin{align}
        A_1^{01} & =\begin{pmatrix}1&1\\0&0\end{pmatrix}.
        \notag
    \end{align}
    \begin{align}
        A_1^{10} & =\begin{pmatrix}0&0\\-1&1\end{pmatrix}.
        \notag
    \end{align}
    Its stacked square is the bond-four tensor
    $\widetilde A_0^{ik}=\sum_jA_1^{ij}\otimes A_1^{jk}$, so that
    $\widetilde A_0^{00}=A_1^{01}\otimes A_1^{10}$,
    $\widetilde A_0^{11}=A_1^{10}\otimes A_1^{01}$, and the two mixed
    matrices vanish~\cite[Section~II]{Liu2025Bialgebraic}. The single
    target of its compression is the bond-one identity tensor
    $\delta^{ij}=\delta_{i,j}$.
\end{definition}

The periodic operator $U_N$ of $A_1$ has the matrix coefficients
$U_N(s,t)=\tr(A_1^{s_0t_0}\cdots A_1^{s_{N-1}t_{N-1}})$, the ordinary
unnormalized trace:
\begin{center}
    % Formula: U_N(s,t)=\tr(A_1^{s_0t_0}\cdots A_1^{s_{N-1}t_{N-1}}), the
    % cyclic trace expanded in the proof of
    % Theorem~\ref{thm:asymex_czy_kernel}.
    % Ink-to-index map: each node is A_1; its upper leg is the output s_n,
    % its lower leg the input t_n, and the horizontal legs are the bond-two
    % virtual index \alpha_n.
    % Contraction set: every horizontal virtual bond, including the return
    % bond that closes the ring.
    % Boundary signature: (virtual-west=0 | virtual-east=0 | physical-up=N |
    % physical-down=N); three sites are drawn and the ellipsis stands for the
    % others.
    \tenkzkernel
    \begin{tenkz}[rows={op}, cols=4, boundary=periodic]
        \tn[skin=mpo, ports={90:physical:$s_0$, 270:physical:$t_0$}]{A_1} &
        \tn[skin=mpo, ports={90:physical:$s_1$, 270:physical:$t_1$}]{A_1} &
        \tn[skin=dots, ports={180:virtual, 0:virtual}]{} &
        \tn[skin=mpo, ports={90:physical:$s_{N-1}$, 270:physical:$t_{N-1}$}]{A_1}
    \end{tenkz}
\end{center}

The periodic
kernel and the identity $U_N^2=I$ are established in
Theorems~\ref{thm:asymex_czy_kernel} and~\ref{thm:asymex_czy_square_identity}.
Theorems~\ref{thm:asymex_czy_compression}
and~\ref{thm:asymex_czy_printed_reduction} give the compression and the
printed bond-three reduction; Theorem~\ref{thm:asymex_czy_reduced_nonsplit}
excludes compressions of the reduced tensor onto the identity with zero
remainder. Arbitrary direct-sum decompositions and the associator mentioned
above are not settled by these results.

\begin{theorem}[Periodic kernel of the CZY operator]
    \label{thm:asymex_czy_kernel}
    \lean{CZYCompression.mpo_czyTensor_apply}
    \leanok
    \uses{def:asymex_czy_tensors}
    For every $N>0$ and $s,t\in\{0,1\}^N$, with site subscripts modulo $N$,
    \begin{align}
        U_N(s,t) & =\delta_{s,\bar t}\prod_{n=0}^{N-1}(-1)^{s_n(1-s_{n+1})},
        \label{eq:asymex_czy_kernel}
    \end{align}
    where $\bar t_n=1-t_n$.
\end{theorem}

\begin{proof}
    \leanok
    The periodic operator is the cyclic trace
    \begin{align}
        U_N(s,t) & =\sum_{\alpha\in\{0,1\}^N}\prod_{n=0}^{N-1}
        (A_1^{s_nt_n})_{\alpha_n\alpha_{n+1}},
        \notag
    \end{align}
    with $\alpha_N=\alpha_0$. The entry $(A_1^{ij})_{\alpha\beta}$ vanishes
    unless $i=1-j$ and $\alpha=i$, and then equals $(-1)^{i(1-\beta)}$. The
    incoming bond at every site is therefore forced to equal the output bit,
    so a product is nonzero only if $s=\bar t$ and $\alpha_n=s_n$ for every
    $n$. The sum then has the single surviving term
    $\prod_n(-1)^{s_n(1-s_{n+1})}$, which is~(\ref{eq:asymex_czy_kernel}).
\end{proof}

\begin{theorem}[Compression of the stacked square onto the identity]
    \label{thm:asymex_czy_compression}
    \lean{CZYCompression.czySquare_compression, CZYCompression.czySquareLeft,
        CZYCompression.czySquareRight, CZYCompression.czySquare_isReduction,
        CZYCompression.czySquare_dim_eq,
        CZYCompression.czySquare_evalWord_remainder_eq_zero,
        CZYCompression.czySquareTarget_isNormal}
    \leanok
    \uses{def:asymex_czy_tensors, thm:asym_multiblock, def:normal}
    The stacked square carries the block triangular form of
    Theorem~\ref{thm:asym_multiblock} for the single normal slot $\delta$
    with $z=3$ zero slots and $4=1+3$. The compression pair is
    $W=(0,1,-1,0)$ and $V=(0,1,0,0)^{\mathsf T}$; it satisfies $WV=1$ and
    $W\widetilde A_0^wV=\delta^w$ for every word $w$. The remainder
    $R^{ij}=\widetilde A_0^{ij}-V\delta^{ij}W$ satisfies
    \begin{align}
        R^{i_1j_1}R^{i_2j_2}R^{i_3j_3}R^{i_4j_4} & =0
        \notag
    \end{align}
    for every choice of four index pairs.
\end{theorem}

\begin{proof}
    \leanok
    \uses{thm:asymex_explicit_gauge}
    With bond coordinates ordered as in the Kronecker product,
    $\widetilde A_0^{00}=e_1u^{\mathsf T}$ and
    $\widetilde A_0^{11}=e_2v^{\mathsf T}$ with $u=(-1,1,-1,1)$ and
    $v=(-1,-1,1,1)$. In the basis
    \begin{align}
        f_0 & =(0,1,1,0)^{\mathsf T},\quad f_1=(0,1,0,0)^{\mathsf T},\quad
        f_2=(1,0,0,0)^{\mathsf T},\quad f_3=(0,0,0,1)^{\mathsf T}
        \notag
    \end{align}
    both letters annihilate $f_0$, have image in $\spn\{f_0,f_1\}$, and
    satisfy $\widetilde A_0^{00}f_1=f_1$ and
    $\widetilde A_0^{11}f_1=f_1-f_0$. Hence every letter is upper
    triangular with diagonal $(0,\delta^{ij},0,0)$. The row $W$ is the
    second row of the inverse of the basis matrix and $V=f_1$. The
    target is normal because its matrix at $(0,0)$ is the scalar $1$.
    The nilpotency is clause~(vi) of Theorem~\ref{thm:asym_multiblock}
    with $|S|+z=4$.
\end{proof}

In a picture, the compression pair closes the bond of the stacked square
onto the bond-one identity tensor, letter by letter:
\begin{center}
    % Formula: W\widetilde A_0^{ij}V=\delta^{ij}, the one-letter case of
    % W\widetilde A_0^wV=\delta^w in Theorem~\ref{thm:asymex_czy_compression}.
    % Ink-to-index map: W is the row vector (0,1,-1,0) and V the column vector
    % (0,1,0,0)^T on the bond \C^4=\C^2\otimes\C^2 of \widetilde A_0; the
    % physical legs are the output i and the input j; \delta is the bond-one
    % tensor \delta^{ij}=\delta_{i,j}.
    % Contraction set: the two virtual bonds joining W and V to
    % \widetilde A_0.
    % Boundary signature: (virtual-west=0 | virtual-east=0 | physical-up=1 |
    % physical-down=1) on both sides.
    \tenkzkernel
    \begin{tenkz}[rows={wire}, cols=3, bonds=none]
        \tn[at=(1,1), name=W, skin=ring, ports={0:virtual}]{W}
        \tn[at=(1,2), name=A, skin=mpo, size=l,
            ports={90:physical:$i$, 270:physical:$j$, 180:virtual, 0:virtual}]{\widetilde A_0}
        \tn[at=(1,3), name=V, skin=ring, ports={180:virtual}]{V}
        \tnwire{W.0}{A.180} \tnwire{A.0}{V.180}
    \end{tenkz}
    \quad $=$ \quad
    \begin{tenkz}[rows={wire}, cols=1, bonds=none]
        \tn[at=(1,1), skin=mpo, ports={90:physical:$i$, 270:physical:$j$}]{\delta}
    \end{tenkz}
\end{center}

The next theorem states that the two stacked copies of $A_1$ contract,
on every ring of $N\geq1$ sites, to the identity:
\begin{center}
    % Formula: U_N^2=\Id (Theorem~\ref{thm:asymex_czy_square_identity}), with
    % entries \sum_j U_N(i,j)U_N(j,k)=\delta_{i,k}.
    % Ink-to-index map: on the left, both rows are A_1, the upper row with
    % output i_n and the lower row with input k_n, and the vertical wire at
    % site n is the summed intermediate label j_n; on the right, the bare
    % vertical wires are the identity \delta_{i_n,k_n} at each site.
    % Contraction set: on the left, every vertical wire between the rows,
    % every horizontal virtual bond of each row, and the periodic return bond
    % of each ring; nothing is contracted on the right.
    % Boundary signature: (virtual-west=0 | virtual-east=0 | physical-up=N |
    % physical-down=N) on both sides; three sites are drawn and the ellipsis
    % stands for the others.
    \tenkzkernel
    \begin{tenkz}[rows={op,op}, cols=4, boundary=periodic, bonds=none]
        \tn[at=(1,1), name=a1, skin=mpo, ports={90:physical:$i_0$, 270:physical, 180:virtual, 0:virtual}]{A_1}
        \tn[at=(1,2), name=a2, skin=mpo, ports={90:physical:$i_1$, 270:physical, 180:virtual, 0:virtual}]{A_1}
        \tn[at=(1,3), name=a3, skin=dots, ports={180:virtual, 0:virtual}]{}
        \tn[at=(1,4), name=a4, skin=mpo, ports={90:physical:$i_{N-1}$, 270:physical, 180:virtual, 0:virtual}]{A_1}
        \tn[at=(2,1), name=b1, skin=mpo, ports={90:physical, 270:physical:$k_0$, 180:virtual, 0:virtual}]{A_1}
        \tn[at=(2,2), name=b2, skin=mpo, ports={90:physical, 270:physical:$k_1$, 180:virtual, 0:virtual}]{A_1}
        \tn[at=(2,3), name=b3, skin=dots, ports={180:virtual, 0:virtual}]{}
        \tn[at=(2,4), name=b4, skin=mpo, ports={90:physical, 270:physical:$k_{N-1}$, 180:virtual, 0:virtual}]{A_1}
        \tnwire{a1.0}{a2.180} \tnwire{a2.0}{a3.180} \tnwire{a3.0}{a4.180}
        \tnwire{b1.0}{b2.180} \tnwire{b2.0}{b3.180} \tnwire{b3.0}{b4.180}
        \tnwire{a1.270}{b1.90} \tnwire{a2.270}{b2.90} \tnwire{a4.270}{b4.90}
    \end{tenkz}
    \quad $=$ \quad
    \begin{tenkz}[rows={wire}, cols=4, bonds=none]
        \tn[at=(1,1), skin=none, ports={90:physical:$i_0$, 270:physical:$k_0$}]{}
        \tn[at=(1,2), skin=none, ports={90:physical:$i_1$, 270:physical:$k_1$}]{}
        \tn[at=(1,3), skin=dots]{}
        \tn[at=(1,4), skin=none, ports={90:physical:$i_{N-1}$, 270:physical:$k_{N-1}$}]{}
    \end{tenkz}
\end{center}

\begin{theorem}[The stacked operator is the identity at every positive length]
    \label{thm:asymex_czy_square_identity}
    \lean{CZYCompression.czySquare_trace_evalWord, CZYCompression.mpo_czySquare,
        CZYCompression.mpo_czyTensor_mul_self,
        CZYCompression.mpo_czySquare_zero_ne_one}
    \leanok
    \uses{thm:asymex_czy_compression, prop:asym_compression_trace,
        lem:asym_mpo_trace_eq, thm:mpdo_operator_product_mpo}
    For every nonempty word $w$, $\tr(\widetilde A_0^w)=\tr(\delta^w)$.
    Consequently $U_N^2=\Id$ for every $N\geq1$, odd $N$ included. At
    $N=0$ the empty periodic contraction of the stacked square is
    $\tr(\Id_4)=4$, so the restriction $N\geq1$ is necessary.
\end{theorem}

\begin{proof}
    \leanok
    Proposition~\ref{prop:asym_compression_trace} applied to the
    compression of Theorem~\ref{thm:asymex_czy_compression} gives the
    trace identity, and Lemma~\ref{lem:asym_mpo_trace_eq} turns it into
    the identity of periodic operators, since $\delta$ generates $\Id$.
    The product formula of the operator of a stacked tensor gives
    $U_N^2=\Id$.
\end{proof}

\begin{theorem}[Nonsplitting of the stacked square]
    \label{thm:asymex_czy_no_intertwiner}
    \lean{CZYCompression.czySquare_right_intertwiner_eq_zero,
        CZYCompression.czySquare_left_intertwiner_eq_zero,
        CZYCompression.czySquare_remainder_ne_zero,
        CZYCompression.czySquare_not_starClosed}
    \leanok
    \uses{thm:asymex_czy_compression, cor:asym_anomaly_not_star}
    If $X\in\C^{4\times1}$ satisfies $\widetilde A_0^{ij}X=X\delta^{ij}$
    for every index pair, then $X=0$; likewise for row vectors $Y$ with
    $Y\widetilde A_0^{ij}=\delta^{ij}Y$. Hence every multi-block
    compression of $\widetilde A_0$ onto $\delta$ has a nonzero remainder,
    and the algebra generated by the matrices of $\widetilde A_0$ is not
    closed under conjugate transposition.
\end{theorem}

\begin{proof}
    \leanok
    From $\widetilde A_0^{00}X=X$ the vector $X$ is a multiple of $e_1$,
    and $\widetilde A_0^{11}e_1=-e_2$ forces the multiple to vanish. For
    rows, $Y\widetilde A_0^{00}=Y$ gives $Y=y_1u^{\mathsf T}$, and
    $Y\widetilde A_0^{11}=Y$ then gives $y_1=0$. A vanishing remainder
    would make $V$ a sitewise intertwiner with $WV=1$; the last assertion
    is Corollary~\ref{cor:asym_anomaly_not_star}.
\end{proof}

\begin{definition}[The reduced tensor]
    \label{def:asymex_czy_reduced}
    \lean{CZYCompression.czyReducedTensor, CZYCompression.czyReduced}
    \leanok
    \uses{def:mpo_tensor, def:mpo_to_mps}
    The tensor $A_0$ of~\cite[Section~II]{Liu2025Bialgebraic} has bond
    dimension three, vanishing mixed matrices, and
    \begin{align}
        A_0^{00} & =\begin{pmatrix}0&-1&1\\0&1&-1\\0&0&0\end{pmatrix}.
        \notag
    \end{align}
    \begin{align}
        A_0^{11} & =\begin{pmatrix}0&1&1\\0&1&1\\0&0&0\end{pmatrix}.
        \notag
    \end{align}
\end{definition}

\begin{theorem}[The printed change of basis]
    \label{thm:asymex_czy_printed_reduction}
    \lean{CZYCompression.czyX11, CZYCompression.czyX11_mul_transpose,
        CZYCompression.czyX11_conj, CZYCompression.czyReducedPadInt_apply}
    \leanok
    \uses{def:asymex_czy_tensors, def:asymex_czy_reduced}
    The matrix
    \begin{align}
        X_{1,1} & =\frac{1}{\sqrt2}\begin{pmatrix}
                                       0  & 1  & 1 & 0  \\
                                       0  & -1 & 1 & 0  \\
                                       -1 & 0  & 0 & 1  \\
                                       -1 & 0  & 0 & -1
                                   \end{pmatrix}
        \notag
    \end{align}
    of~\cite[Appendix~B]{Liu2025Bialgebraic} is orthogonal, and
    $X_{1,1}\widetilde A_0^{ij}X_{1,1}^{-1}=A_0^{ij}\oplus0$ for every
    index pair, exactly as printed.
\end{theorem}

\begin{proof}
    \leanok
    Both identities are integer matrix identities after multiplying by
    $\sqrt2$ on each side: $(\sqrt2X_{1,1})(\sqrt2X_{1,1})^{\mathsf T}=2\Id$
    and
    $(\sqrt2X_{1,1})\widetilde A_0^{ij}(\sqrt2X_{1,1})^{\mathsf T}
        =2(A_0^{ij}\oplus0)$.
\end{proof}

\begin{theorem}[Compression of the reduced tensor]
    \label{thm:asymex_czy_reduced_compression}
    \lean{CZYCompression.czyReduced_compression,
        CZYCompression.czyReduced_evalWord_remainder_eq_zero,
        CZYCompression.mpo_czyReducedTensor}
    \leanok
    \uses{def:asymex_czy_reduced, thm:asym_multiblock, lem:asym_mpo_trace_eq}
    In its own coordinates, $A_0$ carries the block triangular form of
    Theorem~\ref{thm:asym_multiblock} for the single slot $\delta$ with
    $z=2$ zero slots, ordered as a zero slot, the slot $\delta$, and a zero
    slot. The remainder $R^{ij}$ of this compression satisfies
    \begin{align}
        R^{i_1j_1}R^{i_2j_2}R^{i_3j_3} & =0
        \notag
    \end{align}
    for every choice of three index pairs, and the periodic operator of $A_0$
    is the identity at every length $N\geq1$.
\end{theorem}

\begin{proof}
    \leanok
    \uses{thm:asymex_explicit_gauge}
    The matrices $A_0^{00}$ and $A_0^{11}$ are upper triangular with
    diagonal $(0,1,0)$, and the mixed matrices vanish, so the diagonal is
    $(0,\delta^{ij},0)$. The operator identity follows as in
    Theorem~\ref{thm:asymex_czy_square_identity}.
\end{proof}

\begin{theorem}[The reduced tensor does not split]
    \label{thm:asymex_czy_reduced_nonsplit}
    \lean{CZYCompression.czyReduced_right_intertwiner_eq_zero,
        CZYCompression.czyReduced_left_intertwiner_eq_zero,
        CZYCompression.czyReduced_remainder_ne_zero,
        CZYCompression.czyReduced_not_starClosed,
        CZYCompression.czyReduced_adjoin_ne_top}
    \leanok
    \uses{thm:asymex_czy_reduced_compression, cor:asym_anomaly_not_star}
    Both sitewise intertwiner spaces between $A_0$ and $\delta$ vanish.
    Hence every multi-block compression of $A_0$ onto $\delta$ has a
    nonzero remainder, the algebra generated by the matrices of $A_0$ is
    not closed under conjugate transposition, and in particular it is a
    proper subalgebra of the $3\times3$ matrices. In particular $A_0$ is
    not injective, as stated in~\cite[Section~II]{Liu2025Bialgebraic}.
\end{theorem}

\begin{proof}
    \leanok
    For $X=(a,b,c)^{\mathsf T}$, $A_0^{00}X=X$ gives $c=0$ and $a=-b$,
    and $A_0^{11}X=X$ then gives $b=0$. For $Y=(p,q,r)$,
    $YA_0^{00}=Y$ gives $p=0$ and $r=-q$, and $YA_0^{11}=Y$ gives $q=0$.
    The remaining assertions follow as in
    Theorem~\ref{thm:asymex_czy_no_intertwiner}; the full matrix algebra is
    closed under conjugate transposition.
\end{proof}


\begin{theorem}[Indecomposability of the reduced virtual space]
    \label{thm:asymex_czy_reduced_indecomposable}
    \lean{CZYCompression.czyReduced_commutant_form,
        CZYCompression.czyReduced_commuting_idempotent,
        CZYCompression.czyReduced_invariant_isCompl}
    \leanok
    \uses{def:asymex_czy_reduced, thm:asymex_czy_reduced_nonsplit}
    Every matrix commuting with $A_0^{00}$ and $A_0^{11}$ has the form
    $aI+bE_{13}$. The only idempotents in this commutant are $0$ and $I$.
    Consequently, if $\mathbb C^3=V\oplus W$ and both summands are invariant
    under the letters of $A_0$, then $V=0$ or $W=0$.
    In particular, no simultaneous bond similarity decomposes $A_0$ into
    a direct sum of injective matrix product operators, as asserted
    in~\cite[Section~II]{Liu2025Bialgebraic}.
\end{theorem}

\begin{proof}
    \leanok
    Comparing the entries of $PA_0^{00}=A_0^{00}P$ and
    $PA_0^{11}=A_0^{11}P$ gives $P=aI+bE_{13}$.
    Since $E_{13}^2=0$, the equation $P^2=P$ gives $a^2=a$ and $2ab=b$.
    Thus $a\in\{0,1\}$ and $b=0$.
    For an invariant decomposition $V\oplus W$, the projection onto $V$
    along $W$ commutes with every letter and is idempotent. Its range is
    therefore zero or the whole space.
    Any nontrivial direct sum of injective blocks would give such an
    invariant decomposition; a single injective block is excluded by
    Theorem~\ref{thm:asymex_czy_reduced_nonsplit}.
\end{proof}

\begin{theorem}[A one-sided fusion identity]
    \label{thm:asymex_czy_fusion}
    \lean{CZYCompression.czyY00_fusion}
    \leanok
    \uses{def:asymex_czy_reduced}
    Let $Y_{0,0}$ be the first three rows of the matrix $X_{0,0}$
    of~\cite[Appendix~B]{Liu2025Bialgebraic}:
    \begin{align}
        Y_{0,0} & =\begin{pmatrix}
                       0 & 1 & 0 & 1 & 0 & 0 & 0 & 0 & 0 \\
                       0 & 0 & 0 & 0 & 2 & 0 & 0 & 0 & 2 \\
                       0 & 0 & 0 & 0 & 0 & 2 & 0 & 2 & 0
                   \end{pmatrix}.
        \notag
    \end{align}
    Then $Y_{0,0}\sum_jA_0^{ij}\otimes A_0^{jk}=A_0^{ik}Y_{0,0}$ for all
    $i,k$.
\end{theorem}

\begin{proof}
    \leanok
    Direct multiplication of integer matrices.
\end{proof}

The identity of Theorem~\ref{thm:asymex_czy_fusion} moves $Y_{0,0}$ from the
west of two stacked copies of $A_0$ to the east of one copy:
\begin{center}
    % Formula: Y_{0,0}\sum_jA_0^{ij}\otimes A_0^{jk}=A_0^{ik}Y_{0,0}
    % (Theorem~\ref{thm:asymex_czy_fusion}); compare the fusion tensors of
    % \cite[Appendix~B]{Liu2025Bialgebraic}.
    % Ink-to-index map: the tall box is Y_{0,0}, a map from \C^3\otimes\C^3 to
    % \C^3, with its single bond leg on the west and its two product-bond legs
    % on the east. On the left, the upper A_0 carries the output i, the lower
    % A_0 the input k, and the vertical wire between them is the summed label
    % j; on the right, one A_0 carries the output i and the input k.
    % Contraction set: on the left, the vertical wire j and the two
    % product-bond legs of Y_{0,0}; on the right, the bond joining A_0 to the
    % single leg of Y_{0,0}.
    % Boundary signature: (virtual-west=1 | virtual-east=2 | physical-up=1 |
    % physical-down=1) on both sides.
    \tenkzkernel
    \begin{tenkz}[rows={op,op}, cols=2, bonds=none]
        \tnfuse[at=(1,1), name=Y, skin=pill]{Y_{0,0}}
        \tn[at=(1,2), name=u, skin=mpo, ports={90:physical:$i$, 270:physical, 180:virtual, 0:virtual}]{A_0}
        \tn[at=(2,2), name=d, skin=mpo, ports={90:physical, 270:physical:$k$, 180:virtual, 0:virtual}]{A_0}
        \tnwire{open w}{Y.180}
        \tnwire{Y.0@1}{u.180} \tnwire{Y.0@2}{d.180}
        \tnwire{u.270}{d.90}
        \tnwire{u.0}{open e} \tnwire{d.0}{open e}
    \end{tenkz}
    \quad $=$ \quad
    \begin{tenkz}[rows={op,op}, cols=2, bonds=none]
        \tn[at=(1,1), name=a, skin=mpo, ports={90:physical:$i$, 270:physical:$k$, 180:virtual, 0:virtual}]{A_0}
        \tnfuse[at=(1,2), name=Y, skin=pill]{Y_{0,0}}
        \tn[at=(2,1), skin=none]{}
        \tnwire{open w}{a.180} \tnwire{a.0}{Y.180}
        \tnwire{Y.0@1}{open e} \tnwire{Y.0@2}{open e}
    \end{tenkz}
\end{center}


\begin{definition}[Printed fusion maps for the nontrivial element]
    \label{def:asymex_czy_printed_fusions}
    \lean{CZYCompression.czyY11, CZYCompression.czyY01, CZYCompression.czyY10,
        CZYCompression.czyY11_eq_submatrix}
    \leanok
    \uses{def:asymex_czy_reduced, thm:asymex_czy_printed_reduction}
    As in~\cite[Appendix~B]{Liu2025Bialgebraic}, $Y_{1,1}$ is the first
    three rows of $X_{1,1}$, and
    \begin{align}
        Y_{0,1} & =\begin{pmatrix}0&0&1&0&0&-1\\0&0&0&1&-1&0\end{pmatrix}, &
        Y_{1,0} & =\begin{pmatrix}0&1&0&0&0&1\\0&0&1&0&1&0\end{pmatrix}.
        \notag
    \end{align}
\end{definition}

\begin{theorem}[Local fusion and the nontrivial triple]
    \label{thm:asymex_czy_nontrivial_associator}
    \lean{CZYCompression.czyY01_fusion, CZYCompression.czyY10_fusion,
        CZYCompression.czyY11_fusion, CZYCompression.czyFusion_nontrivial_triple}
    \leanok
    \uses{def:asymex_czy_printed_fusions}
    The printed maps satisfy
    $Y_{a,b}\sum_j A_a^{ij}\otimes A_b^{jk}=A_{a+b}^{ik}Y_{a,b}$
    for $(a,b)=(0,1),(1,0),(1,1)$.
    At the nontrivial triple,
    $Y_{0,1}(Y_{1,1}\otimes I_2)=-Y_{1,0}(I_2\otimes Y_{1,1})$,
    with the normalization of~\cite[Section~II]{Liu2025Bialgebraic}.
\end{theorem}

\begin{proof}
    \leanok
    Substitute the displayed matrices and the letters of $A_0,A_1$.
    Each equality is a finite matrix identity; the common factor
    $1/\sqrt2$ cancels from the triple comparison.
\end{proof}

\begin{theorem}[Associativity on the physical images]
    \label{thm:asymex_czy_supported_associativity}
    \lean{CZYCompression.czyDim, CZYCompression.czyFamilyInt,
        CZYCompression.czyFusionInt, CZYCompression.czyTripleInt,
        CZYCompression.czyLeftInt, CZYCompression.czyRightInt,
        CZYCompression.czyFusion_supported_associativity_int,
        CZYCompression.czyFusion_supported_associativity,
        CZYCompression.czyFusion_supported_associativity_word,
        CZYCompression.czyFusion_printed_neutral_no_scalar}
    \leanok
    \uses{thm:asymex_czy_nontrivial_associator, thm:asymex_czy_fusion}
    Put $F_{0,0}=Y_{0,0}$, $F_{0,1}=2Y_{0,1}$,
    $F_{1,0}=2Y_{1,0}$ and $F_{1,1}=2\sqrt2Y_{1,1}$.
    For $a,b,c\in\mathbb Z_2$, define
    $T^{ik}=\sum_{j,l}A_a^{ij}\otimes A_b^{jl}\otimes A_c^{lk}$,
    $L=F_{a+b,c}(F_{a,b}\otimes I_c)$ and
    $R=F_{a,b+c}(I_a\otimes F_{b,c})$, identifying both bond orders.
    Then $LT^{ik}=(-1)^{abc}RT^{ik}$ for every physical pair $(i,k)$.
    More generally, $LT^{p_1}\cdots T^{p_n}=(-1)^{abc}RT^{p_1}\cdots T^{p_n}$
    for every nonempty physical word.
    This is the corrected relation on the physical images
    of~\cite{gap:lmsvkl25_czy_associator}.
    On the full virtual space the printed neutral fusions admit no
    scalar comparison: the entry in row one and column $(3,3,1)$ is two
    on the left and zero on the right.
\end{theorem}

\begin{proof}
    \leanok
    For the chosen normalization all matrices are integral. Direct
    multiplication proves the identity for the eight group triples and
    the four physical pairs. Multiplication on the right by the
    remaining letters gives the nonempty-word identity.
    The specified entry proves the full-space obstruction.
\end{proof}
